\chapter{Аннотация}
\begin{SingleSpace}

\noindent Ожоговая травма - это тяжелое и общемировое бремя здравоохранения - по данным Всемирной организации здравоохранения, каждый год от ожогов страдает порядка 11 миллионов человек. Правильная и своевременная классификация степени ожога является критически важной для адекватного выбора лечебной тактики: ожоги 1 степени требуют амбулаторного лечения, ожоги 3 степени требуют госпитализации с возможным хирургическим вмешательством.

\noindent В данной диссертации разработана система автоматической классификации степеней ожогов на основе глубокого обучения, использующая архитектуру YOLOv8. Модель обучалась на наборе данных Kaggle Skin Burn Dataset, содержащем 1\,225 изображений ожоговых поражений кожи и распределенных по трем классам: ожоги 1 степени (поверхностные), 2 степени (частичной толщины) и 3 степени (полной толщины).

\noindent Предложенная модель на базе YOLOv8n (nano) показала mAP@0.5 = 0.847, precision = 0.823, recall = 0.791. Архитектура YOLOv8 была выбрана за возможность детектирования объектов в реальном времени и высокую производительность для медицинских изображений. Обучение модели проводилось на облачной GPU платформе vast.ai с использованием трансферного обучения на предобученные веса COCO.

\noindent Научная новизна заключается в использовании новой архитектуры детекции объектов YOLOv8 для задачи классификации ожоговых поражений, в формировании конвейера для обработки медицинских изображений с аугментацией данных и в разработке REST API для практического использования модели в клинических условиях.

\vspace{0.8em}
\keywords{ожоговая болезнь, глубокое обучение, YOLOv8, детекция объектов, медицинская визуализация, классификация изображений, компьютерное зрение}

\end{SingleSpace}